\subsection{The Euclidean selection inequality}
\label{geo:angles}

This subsection proves \eqref{geo:four-selection} and records all
constants used in the contact theorem. Put
\[
 f_t(\theta)=t\sin\theta-1-\cos\theta,\qquad
 g_t(\theta)=4t\sin\theta-2(1-t^2)(1-\cos\theta),\qquad
 H_t=g_t+4f_t.
\]

\paragraph{Planar balanced triples.}
Let three unit vectors in a plane enclose the origin, with cyclic gaps
\(\alpha,\beta,\gamma\in[0,\pi]\) summing to \(2\pi\).
For pair \(i,j\) and singleton \(k\), define \(w,v,A\) as in
\eqref{geo:pair-cell}, and set
\[
 D_k=|w|^2+t^2|v|^2+2t|w\times v|-A^2.
\]
All cyclic oriented sine terms are nonnegative. Expanding gives
\begin{equation}
 D_k=t^2(a_i-a_j)^2+a_ia_jg_t(\theta_{ij})
 +4a_k\{a_if_t(\theta_{ki})+a_jf_t(\theta_{jk})\}.
 \label{geo:planar-deficit}
\end{equation}
Take \(t=3/4\), \(\theta_0=2\arctan(1/t)\), and
\(\theta_1=2\arctan t=\pi-\theta_0\).
One has \(f_t\ge0\) on \([\theta_0,\pi]\) and
\(g_t\ge0\) on \([0,4\arctan t]\).
If two gaps are at least \(\theta_0\), choose their common vertex as
singleton; the remaining gap is at most \(4\arctan t\), so every
term in \eqref{geo:planar-deficit} is nonnegative.

Otherwise exactly one gap \(\alpha\) is at least \(\theta_0\), as
\(\theta_0<2\pi/3\). Its other gaps lie in
\([\theta_1,\theta_0]\). Add the deficits at the two endpoints \(i,j\)
of \(\alpha\):
\[
 D_i+D_j=t^2\{(a_j-a_k)^2+(a_i-a_k)^2\}
 +8a_ia_jf_t(\alpha)+a_ia_kH_t(\gamma)+a_ja_kH_t(\beta).
\]
On that middle interval,
\(\sin\theta\ge2t/(1+t^2)\) and
\(\cos\theta\le(1-t^2)/(1+t^2)\), hence
\[
 H_t(\theta)\ge
 \frac{4(t^4+3t^2-2)}{1+t^2}=\frac1{100}.
\]
Thus some deficit is nonnegative. Monotonicity in \(t\) gives
\eqref{geo:four-selection} for every balanced subset of at most
three rays, including limiting degeneracies.

\paragraph{A nearest-facet inequality.}
Now suppose four rays form a tetrahedron with the origin in its
interior; all other cases have a balanced subset of at most three.
Choose a facet of least plane distance \(h>0\) from the origin,
with outward normal \(e\), and write its vertices
\[
 n_i=he+Ru_i,\qquad R=\sqrt{1-h^2},\qquad u_i\in S^1\subset e^\perp.
\]
The radius-\(h\) ball lies in the tetrahedron and touches this facet
in its relative interior: equality with another facet would force
its unit normal to equal \(e\).
The projected triangle therefore encloses zero. Its cyclic gaps
\(\alpha,\beta,\gamma\) are in \((0,\pi)\), and
\[
 \Pi=\cos(\alpha/2)\cos(\beta/2)\cos(\gamma/2)>0,\qquad
 p=\frac{h^2}{R^2}.
\]
We claim
\begin{equation}
 \frac{h}{1-h}\le\sqrt{2\Pi},\qquad
 p\le\frac{2\Pi}{1+2\sqrt{2\Pi}}.
 \label{geo:facet}
\end{equation}
Write the fourth vertex as \(r-se\), \(0<s\le1\).
For an edge let \(\nu\) be its outward planar unit normal,
\(b=R\cos(\theta/2)>0\) its chord distance, and
\(\ell=2R\sin(\theta/2)\) its length. The adjacent facet equation is
\[
 (s+h)\nu\cdot x+(\nu\cdot r-b)z=sb+h\nu\cdot r.
\]
Its right side is positive, since the origin is in the interior.
Its distance from the origin is at least \(h\); squaring yields
\[
 \nu\cdot r\ge
 \frac{h(s+h)}{2b}-\frac{(s-h)b}{2h}.
\]
Multiply by \(\ell\), sum, and use the polygon identity
\(\sum\ell\nu=0\). This forces \(s>h\) and gives
\[
 h^2\frac{s+h}{s-h}\le\frac{\sum\ell b}{\sum\ell/b}
 =2R^2\Pi.
\]
The equality follows from the half-gap sum \(\pi\) and the identities
\(\sum\sin2A=4\sin A\sin B\sin C\),
\(\sum\tan A=\tan A\tan B\tan C\).
Since \(s\le1\), this gives \(h^2\le2(1-h)^2\Pi\), proving
\eqref{geo:facet}. In particular, concavity of \(\log\cos\) gives
\(\Pi\le1/8\) and \(h\le1/3\).

\paragraph{Lifting the planar deficit.}
Set \(t=9/10\). If \(w_0,v_0\) are the vectors for the projected
triangle, the normal component of \(w\times v\) has magnitude
\(R^2|w_0\times v_0|\), while
\[
 w=Rw_0+h(a_i+a_j-2a_k)e,\qquad
 v=Rv_0+h(a_i-a_j)e.
\]
Consequently
\[
 D_k\ge R^2D_k^{\rm planar}
 +h^2\{t^2(a_i-a_j)^2-8a_k(a_i+a_j)\}.
\]
Dropping the last positive square and using
\eqref{geo:planar-deficit},
\begin{equation}
 \frac{D_k}{R^2}\ge
 t^2(a_i-a_j)^2+a_ia_jg_t(\theta_{ij})
 +4a_k\{a_i(f_t(\theta_{ki})-2p)
             +a_j(f_t(\theta_{jk})-2p)\}.
 \label{geo:lifted}
\end{equation}
Call a gap large when it is at least $2\pi/3$. We prove the following
three bounds from \eqref{geo:facet}:
\begin{align}
 f_t(\theta)-2p&\ge0
 &&\text{for either large gap if two are large},
 \label{geo:coeff1}\\
 f_t(\alpha)-2p&\ge-1/200
 &&\text{for every }\alpha\ge2\pi/3,
 \label{geo:coeff2}\\
 H_t(\beta)-8p,\ H_t(\gamma)-8p&\ge1/25
 &&\text{if }\alpha\text{ is the unique large gap}.
 \label{geo:coeff3}
\end{align}

For \eqref{geo:coeff1}, put \(B=\theta/2\in[\pi/3,\pi/2]\),
\(z=\cos B\), and
\[
 r=\frac{\sqrt3}{2}\sin B-\frac z2,\qquad
 A=\frac{\sqrt3}{2}\sin B-z.
\]
The second large gap implies \(2\Pi\le zr\).
Both \(r,A\) increase with \(B\), while \(z\) decreases.
At \(B=\pi/3\), \(16rA^2=z\), so \(16rA^2\ge z\) throughout.
By \eqref{geo:facet} and monotonicity of \(u/(1+2\sqrt u)\),
\[
 p\le\frac{zr}{1+2\sqrt{zr}}\le zA
 =\frac12 f_{\sqrt3/2}(2B).
\]
The middle inequality is equivalent, by positive squaring, to the
inequality just proved. Since \(9/10>\sqrt3/2\),
\eqref{geo:coeff1} follows.

For \eqref{geo:coeff2}, put \(z=\cos(\alpha/2)\in[0,1/2]\).
The product of the other two half-gap cosines is at most
\((1-z)/2\). Thus
\[
 f_t(\alpha)-2p
 \ge2z\left[\frac9{10}\sqrt{1-z^2}-q(z)\right],\qquad
 q(z)=z+\frac{1-z}{1+2\sqrt{z(1-z)}}.
\]
On \(0\le z\le1/100\), use \(q(z)\le1\) and
\(\sqrt{1-z^2}\ge1-z^2\), obtaining
\[
 f_t(\alpha)-2p
 \ge-\frac{10009}{5000000}>-\frac1{200}.
\]
On \(1/100\le z\le1/4\), let \(k=\sqrt{z(1-z)}\).
The sign of \(q'\) is the sign of
\[
 N=(1-z)(4z-1)+4k^3,\qquad N'=5-8z+6k(1-2z)>0.
\]
Hence \(q\) has no interior maximum. Its endpoint values are at
most \(1/100+99/118<17/20\) and \(2/3<17/20\), whereas
\((9/10)\sqrt{1-z^2}\ge(9/10)\sqrt{15/16}>17/20\).
On \(1/4\le z\le1/2\), \(N\ge0\), so \(q(z)\le q(1/2)=3/4\);
meanwhile \((9/10)\sqrt{1-z^2}\ge9\sqrt3/20>3/4\).
These bounds prove \eqref{geo:coeff2}.

For \eqref{geo:coeff3}, the two middle gaps belong to
\([\pi/3,2\pi/3]\). For either such \(\theta\), maximizing the other
two half-gap cosine factors subject to one being large gives
\[
 p\le2\Pi\le\frac{\sqrt3\sin\theta-1-\cos\theta}{4}.
\]
Since
\(H_{9/10}=(36/5)\sin\theta-219/50-(181/50)\cos\theta\),
\[
 H_{9/10}(\theta)-8p\ge
 \left(\frac{36}{5}-2\sqrt3\right)\sin\theta
             -\frac{119}{50}-\frac{81}{50}\cos\theta.
\]
This expression is concave on the stated interval. Its smaller
endpoint value is \(18\sqrt3/5-619/100>1/25\), the last inequality
following after positive squaring from \(388800>388129\).

\paragraph{Completing selection.}
If two gaps are at least \(2\pi/3\), choose their common vertex
as singleton. The remaining gap is at most \(2\pi/3\), where
\(g_{9/10}\ge0\); \eqref{geo:coeff1} makes \eqref{geo:lifted}
nonnegative.
Otherwise let \(a,b\) be the amplitudes at the endpoints of the
unique large gap, and \(c\) the third amplitude. Add the endpoint
deficits, and use \eqref{geo:coeff2}--\eqref{geo:coeff3}:
\[
 \frac{D_i+D_j}{R^2}\ge
 \frac{81}{100}\{(a-c)^2+(b-c)^2\}
      -\frac{ab}{25}+\frac{c(a+b)}{25}
 =\frac{79}{200}(a+b-2c)^2
  +\frac{83}{200}(a-b)^2+\frac{c^2}{25}\ge0.
\]
Some deficit is nonnegative, proving \eqref{geo:four-selection}.
All amplitudes may vanish; degeneracies are covered either by
continuity or reduction to a smaller balanced subset.

\paragraph{The optional explicit cusp constant.}
For \(d=3\), the shifted diagonal Gram matrix on \(R_1\) is
\[
 M=\frac1{3888}
 \begin{pmatrix}251&-13&-13\\-13&20&-7\\-13&-7&20\end{pmatrix}.
\]
The exact Haar integrals give this matrix; off-diagonal phase
averages have larger unmultiplied coefficients, at least \(7/648\).
For its positive eigenvalues,
\(\lambda_{\min}(M)\ge4\det M/(\operatorname{tr}M)^2=325/338724\).
Also
\(\|z^*Dz-(7/36)\operatorname{tr}D\|_\infty\le(19/12)\|D\|\).
Dividing the \(L^2\) lower bound by this \(L^\infty\) upper bound
gives \(\kappa_3=325/536313\) in Lemma~\ref{geo:cusp}.
